% Official source: https://www.imo-official.org/assets/documents/problems/2023/2023_eng.pdf
\competitionyear{2023}

\begin{problem}{2023}{1}
Determine all composite integers $n>1$ that satisfy the following property: if $d_1,d_2,\dots,d_k$ are all the positive divisors of $n$ with $1=d_1<d_2<\cdots<d_k=n$, then $d_i$ divides $d_{i+1}+d_{i+2}$ for every $1\le i\le k-2$.
\end{problem}

\begin{problem}{2023}{2}
Let $ABC$ be an acute-angled triangle with $AB<AC$. Let $\Omega$ be the circumcircle of $ABC$. Let $S$ be the midpoint of the arc $CB$ of $\Omega$ containing $A$. The perpendicular from $A$ to $BC$ meets $BS$ at $D$ and meets $\Omega$ again at $E\ne A$. The line through $D$ parallel to $BC$ meets line $BE$ at $L$. Denote the circumcircle of triangle $BDL$ by $\omega$. Let $\omega$ meet $\Omega$ again at $P\ne B$.

Prove that the line tangent to $\omega$ at $P$ meets line $BS$ on the internal angle bisector of $\angle BAC$.
\end{problem}

\begin{problem}{2023}{3}
For each integer $k\ge2$, determine all infinite sequences of positive integers $a_1,a_2,\dots$ for which there exists a polynomial $P$ of the form $P(x)=x^k+c_{k-1}x^{k-1}+\cdots+c_1x+c_0$, where $c_0,c_1,\dots,c_{k-1}$ are non-negative integers, such that
\[
    P(a_n)=a_{n+1}a_{n+2}\cdots a_{n+k}
\]
for every integer $n\ge1$.
\end{problem}

\begin{problem}{2023}{4}
Let $x_1,x_2,\dots,x_{2023}$ be pairwise different positive real numbers such that
\[
    a_n=\sqrt{(x_1+x_2+\cdots+x_n)
    \left(\frac1{x_1}+\frac1{x_2}+\cdots+\frac1{x_n}\right)}
\]
is an integer for every $n=1,2,\dots,2023$. Prove that $a_{2023}\ge3034$.
\end{problem}

\begin{problem}{2023}{5}
Let $n$ be a positive integer. A \emph{Japanese triangle} consists of $1+2+\cdots+n$ circles arranged in an equilateral triangular shape such that for each $i=1,2,\dots,n$, the $i$-th row contains exactly $i$ circles, exactly one of which is coloured red. A \emph{ninja path} in a Japanese triangle is a sequence of $n$ circles obtained by starting in the top row, then repeatedly going from a circle to one of the two circles immediately below it and finishing in the bottom row. Here is an example of a Japanese triangle with $n=6$, along with a ninja path in that triangle containing two red circles.

% Faithful textbook-style TikZ reconstruction of the official 2023 IMO Problem 5 example.
\begin{center}
\begin{tikzpicture}[x=1cm,y=1cm,line cap=round,line join=round]
    \def\circleRadius{0.235}
    \foreach \row in {1,...,6}{
        \foreach \column in {1,...,\row}{
            \coordinate (r\row c\column) at ({(\column-(\row+1)/2)*0.50},{(6-\row)*0.44});
        }
    }
    \foreach \row in {1,...,6}{
        \foreach \column in {1,...,\row}{
            \draw[line width=0.45pt,fill=white] (r\row c\column) circle (\circleRadius);
        }
    }
    % Exactly one red circle in each row, as in the official example.
    \foreach \r/\c in {1/1,2/1,3/3,4/3,5/4,6/1}{
        \draw[line width=0.45pt,fill=exampleRed] (r\r c\c) circle (\circleRadius);
    }
    % A ninja path through two red circles (rows 1 and 4).
    \draw[line width=1.15pt] (r1c1.center)--(r2c2.center)--(r3c2.center)
        --(r4c3.center)--(r5c3.center)--(r6c3.center);
    \draw[latex-latex,line width=0.45pt] (-1.50,-0.37)--(1.50,-0.37);
    \node[below] at (0,-0.39) {$n=6$};
\end{tikzpicture}
\end{center}


In terms of $n$, find the greatest $k$ such that in each Japanese triangle there is a ninja path containing at least $k$ red circles.
\end{problem}

\begin{problem}{2023}{6}
Let $ABC$ be an equilateral triangle. Let $A_1,B_1,C_1$ be interior points of $ABC$ such that $BA_1=A_1C$, $CB_1=B_1A$, $AC_1=C_1B$, and
\[
    \angle BA_1C+\angle CB_1A+\angle AC_1B=480^\circ.
\]
Let $BC_1$ and $CB_1$ meet at $A_2$, let $CA_1$ and $AC_1$ meet at $B_2$, and let $AB_1$ and $BA_1$ meet at $C_2$. Prove that if triangle $A_1B_1C_1$ is scalene, then the three circumcircles of triangles $AA_1A_2$, $BB_1B_2$ and $CC_1C_2$ all pass through two common points.

(Note: a scalene triangle is one where no two sides have equal length.)
\end{problem}
